\section{Exact pivot counts for the algebraic screening family}
\label{sec:pivot-appendix}

The oracle-call separation in Theorem~\ref{thm:lp-family} can be strengthened
for the exact simplex implementation supplied with the repository. This
appendix proves the pivot counts, including the entering and leaving choices;
it does not infer them from the experimental table. The statement concerns
the algebraic family~\eqref{eq:lp-family}, with its specified ordering and
solver. No realization by three-manifold triangulations is asserted.

\begin{theorem}[Exact pivot separation]\label{thm:lp-pivots}
For the $s$-block algebraic family of Theorem~\ref{thm:lp-family}, run the
specified all-slack Bland solver without a pivot cutoff, or with an allowance
sufficient for both complete propagation runs. The total numbers of simplex
pivots are
\begin{equation}\label{eq:lp-pivot-counts}
 P_{\rm base}(s)=6s^3+3s^2+s,
 \qquad
 P_{\rm screened}(s)=3s^2+s.
\end{equation}
Every positive current-node or irrelevant-coordinate deletion query uses
exactly $2s$ pivots. Deleting the mandatory type-zero coordinate of block
$j$, where $0\le j<s$, uses exactly $2j$ pivots. Global cache capacity zero
suffices for the screened formula.
\end{theorem}

We first establish a tableau lemma that includes all the restrictions
arising during these propagation runs.

\subsection{Ordered chain systems and the exact initial tableau}

Let $m\ge1$. Consider disjoint groups $G_0,\ldots,G_{m-1}$ of nonnegative
original variables, none containing a distinguished variable $z$. A group
may contain several variables with identical columns, or it may be empty.
Write
\[
 X_i=\sum_{x\in G_i}x.
\]
Let $\mathcal Y$ contain all other original variables, whose matching
columns and objective coefficients are zero, and put
$Y=\sum_{y\in\mathcal Y}y$. The matching equations, in this exact order,
and the objective are
\begin{equation}\label{eq:pivot-chain}
 X_i-z=0\quad(0\le i<m),\qquad \text{objective }z.
\end{equation}
The normalization used by the solver is
\begin{equation}\label{eq:pivot-normalization}
 z+\sum_{i=0}^{m-1}X_i+Y\le1.
\end{equation}

The implementation forms the inequalities $Cx\le0$, followed by
$-Cx\le0$, followed by~\eqref{eq:pivot-normalization}. It adds one
nonnegative slack to each row and starts with the all-slack feasible basis.
Denote the slacks of the $C$ rows by $a_0,\ldots,a_{m-1}$, those of the
$-C$ rows by $b_0,\ldots,b_{m-1}$, and the normalization slack by $\sigma$.
Thus the initial tableau equations are
\begin{align}
 a_i+X_i-z&=0 &&(0\le i<m),\label{eq:pivot-initial-plus}\\
 b_i-X_i+z&=0 &&(0\le i<m),\label{eq:pivot-initial-minus}\\
 \sigma+z+\sum_{i=0}^{m-1}X_i+Y&=1.\label{eq:pivot-initial-sum}
\end{align}
Every original variable precedes every slack in the variable order.
Slack indices then follow the displayed row order: all $a_i$, all $b_i$,
then $\sigma$. Removing an original column preserves the relative order
of the remaining original variables and of all the slacks.

The entering rule chooses the least-index variable with strictly positive
reduced objective coefficient. The leaving rule minimizes the usual
nonnegative ratio and, among equal ratios, chooses the least-index current
basic variable. Zero reduced coefficients are never selected. The solver
returns immediately once the objective is positive, or when no positive
reduced coefficient remains.

\begin{lemma}[Exact chain pivot path]\label{lem:pivot-chain}
Assume $z$ is present. If every $G_i$ is nonempty, the solver for
\eqref{eq:pivot-chain} returns a positive point after exactly $m+1$ pivots.
If $h$ is the first index for which $G_h$ is empty, it returns a nonpositive
certificate after exactly $h+1$ pivots. These assertions are independent
of the number of irrelevant variables and of the number of available
duplicate columns in each nonempty group.
\end{lemma}
\begin{proof}
Initially $z$ is the only variable with positive reduced coefficient.
Its coefficients are $-1$ in the $C$ rows, $+1$ in the $-C$ rows, and
$+1$ in the normalization row. All $-C$ rows have ratio zero, whereas
the normalization row has ratio one. The leaving variable is therefore
$b_0$, by the specified slack ordering and Bland tie rule. The pivot
coefficient is one, the objective remains zero, and the resulting
dictionary expresses
\[
 z=X_0-b_0.
\]

Here is the full inductive invariant. After $h+1$ pivots, suppose that
$G_0,\ldots,G_{h-1}$ were nonempty. Let $x_i$ be the least-index available
member of $G_i$ that entered for $i<h$, and put $Y_i=X_i-x_i$.
The basic original variables are $z,x_0,\ldots,x_{h-1}$. All $a_i$, all
$b_j$ with $j>h$, and $\sigma$ are also basic. The nonbasic slacks are
$b_0,\ldots,b_h$. The tableau equations are exactly
\begin{align}
 z-X_h+b_h&=0,\label{eq:pivot-invariant-z}\\
 x_i+Y_i-X_h+b_h-b_i&=0 &&(i<h),\label{eq:pivot-invariant-old}\\
 a_i+b_i&=0 &&(i\le h),\label{eq:pivot-invariant-completed}\\
 a_j+X_j-X_h+b_h&=0 &&(j>h),\label{eq:pivot-invariant-plus}\\
 b_j-X_j+X_h-b_h&=0 &&(j>h),\label{eq:pivot-invariant-minus}\\
 \sigma+(h+2)X_h-(h+1)b_h
       +\sum_{i<h}b_i+\sum_{j>h}X_j+Y&=1.
       \label{eq:pivot-invariant-normalization}
\end{align}
The current objective value is zero, and its reduced expression is
\begin{equation}\label{eq:pivot-reduced}
 z=X_h-b_h.
\end{equation}
All tableau right-hand sides except that of the normalization row are
zero. Direct elimination in~\eqref{eq:pivot-initial-plus}--
\eqref{eq:pivot-initial-sum} proves the invariant for $h=0$.

If $G_h$ is empty, equation~\eqref{eq:pivot-reduced} has no positive
reduced coefficient: $b_h$ has coefficient $-1$, and all other nonbasic
variables have coefficient zero. The solver therefore stops after
$h+1$ pivots with its nonpositive certificate.

Suppose instead that $G_h$ is nonempty. Precisely its variables have
positive reduced coefficient, each equal to one. Bland consequently
selects $x_h=\min G_h$. Its coefficient is $-1$ in the $z$ row and in
each earlier $x_i$ row, zero in the $a_i$ rows with $i\le h$, $-1$ in
the later $a_j$ rows, $+1$ in the later $b_j$ rows, and $h+2$ in the
normalization row. Thus the only eligible leaving rows are the later
$b_j$ rows, with ratio zero, and normalization, with ratio $1/(h+2)$.

If $h<m-1$, a zero-ratio row exists. The remaining basic variables
$b_j$, $j>h$, retain their original indices, so Bland chooses $b_{h+1}$.
This is a unit pivot in a row with zero right-hand side. Substitution
using
\[
 X_h=X_{h+1}-b_{h+1}+b_h
\]
gives every invariant equation with $h$ replaced by $h+1$; the new basic
variable $x_h$ supplies the additional instance of
\eqref{eq:pivot-invariant-old}. The normalization coefficients become
$h+3$ and $-(h+2)$, as required. The objective is still zero. This
establishes the induction, including both Bland choices.

If $h=m-1$, there is no later $b_j$ row. The only eligible leaving row
is normalization, with pivot coefficient $m+1$. This final pivot sets
$x_{m-1}=1/(m+1)$ and makes the objective $1/(m+1)>0$.
The implementation returns immediately. The initial $z$ pivot and the
$m$ entering-group pivots give exactly $m+1$ pivots.

Throughout the zero-objective stages, irrelevant variables have zero
reduced coefficient by~\eqref{eq:pivot-reduced}, so none enters.
If a group has several available twins, all have the same positive
coefficient when that group is current and Bland selects its least member.
Afterwards its other twins have reduced coefficient zero. If an earlier
twin was already forbidden before the LP call, the least remaining
member is selected instead. These observations prove the asserted
independence from irrelevant columns and available multiplicities.
\end{proof}

\begin{lemma}[Deletion of the objective variable]\label{lem:pivot-zero-objective}
If the column $z$ is deleted from the restricted query, the solver returns
its nonpositive certificate with zero pivots.
\end{lemma}
\begin{proof}
All remaining original objective coefficients are zero, as are the initial
slack coefficients. The all-slack objective value is zero and there is no
entering variable. The stopping test precedes any pivot.
\end{proof}

\subsection{The actual block ordering and previous propagations}

For the family~\eqref{eq:lp-family}, put $m=2s-1$. The exact matching-row
order used by the model builder is
\begin{align*}
 (Cx)_0&=q_{0,1}+p_0-z,\\
 (Cx)_{2j-1}&=q_{j,0}-z &&(1\le j<s),\\
 (Cx)_{2j}&=q_{j,1}+p_j-z &&(1\le j<s).
\end{align*}
Equivalently, the chain groups before propagation are
\begin{align*}
 G_0&=\{q_{0,1},p_0\},\\
 G_{2j-1}&=\{q_{j,0}\} &&(1\le j<s),\\
 G_{2j}&=\{q_{j,1},p_j\} &&(1\le j<s).
\end{align*}
For completeness, let $n_s=7(2s+\lceil s/4\rceil)$ and index the original
columns $v_0,\ldots,v_{n_s-1}$ in increasing order. In each seven-coordinate
block, the four triangle coordinates precede the three quadrilateral
coordinates. With zero-based indices, the distinguished columns are
\begin{align*}
 q_{j,0}&=v_{7(s+j)+4},&
 q_{j,1}&=v_{7(s+j)+5} &&(0\le j<s),\\
 p_j&=v_{14s+7\lfloor j/4\rfloor+(j\bmod4)}
 &&&&(0\le j<s).
\end{align*}
Thus $z=v_{7s+4}$, and every $q_{j,1}$ precedes its auxiliary triangle
twin $p_j$. These are the block and column conventions of
Section~\ref{sec:witness}. All remaining variables are irrelevant
variables of the chain lemma. The synthetic model has no triangle-anchor
restrictions.

At propagation stage $j$, each earlier block $\ell<j$ has been forced to
quadrilateral type zero. Its variable $q_{\ell,1}$ has therefore been
removed, while $p_\ell$ remains available. Its group $G_{2\ell}$ becomes
the nonempty singleton $\{p_\ell\}$. All other row groups remain nonempty.
Lemma~\ref{lem:pivot-chain} consequently gives exactly
$m+1=2s$ pivots for every current-node LP.

Deleting a tested irrelevant prefix coordinate removes only one member of
$\mathcal Y$. It neither changes the objective nor empties a row group,
so every such deletion query also takes exactly $2s$ pivots. Although the
number of restricted original columns changes, the relative slack order
used in the leaving-variable proof does not.

Deleting the mandatory coordinate $q_{j,0}$ has two different cases.
For $j=0$, this is deletion of $z$, so
Lemma~\ref{lem:pivot-zero-objective} gives zero pivots. For $j>0$, it
empties $G_{2j-1}$ and leaves all earlier groups nonempty. The first empty
row is thus $h=2j-1$, and Lemma~\ref{lem:pivot-chain} gives
$h+1=2j$ pivots.

\subsection{Summing the exact costs}

\begin{proof}[Proof of Theorem~\ref{thm:lp-pivots}]
Theorem~\ref{thm:lp-family} establishes that both runs make $s+1$
current-node queries and $s$ mandatory deletion queries. The former
each cost $2s$ pivots, and the latter cost $0,2,\ldots,2(s-1)$.
The baseline additionally makes $3s^2$ positive irrelevant-prefix
deletion queries, each costing $2s$ pivots. Hence
\begin{align*}
 P_{\rm base}(s)
   &=(s+1)(2s)+(3s^2)(2s)+\sum_{j=0}^{s-1}2j
     =6s^3+3s^2+s,\\
 P_{\rm screened}(s)
   &=(s+1)(2s)+\sum_{j=0}^{s-1}2j
     =3s^2+s.
\end{align*}
Screening these particular deletions needs only the current-node witness,
so a zero-capacity persistent cache gives the same result.
\end{proof}

The difference is exactly $6s^3$ pivots. The proof concerns the specified
matrix family, column and row orders, strict entering rule, and stopping
tests. It establishes neither a polynomial pivot bound for arbitrary
Bland-simplex inputs nor a universal wall-clock factor. It also leaves
the separate geometric realizability question unchanged.
